\documentclass[10pt,letterpaper]{article}
\usepackage{cornell-notes}

\title{Clause 05.09: Preprocessing Numbers}
\author{}
\date{}
\setDocTitle{Clause 05.09: Preprocessing Numbers}
\setDocSubtitle{C++ 2024}
\setDocOwner{C++ 2024 Cornell Notes}
\setDocDate{\today}
\setCornellCollection{C++ 2024 Cornell Notes}
\setCornellUnitType{Clause}
\setCornellUnitNumber{05.09}
\setCornellUnitTitle{Preprocessing Numbers}
\hypersetup{pdftitle={Clause 05.09: Preprocessing Numbers},pdfsubject={Cornell notes},pdfauthor={}}

\begin{document}
\maketitle
\begin{CornellOverviewBox}{Overview}
\textbf{Classification:} Core language - Normative\\[2pt]
\textbf{Learning objective:} Understand the rules, terminology, and semantic consequences associated with Preprocessing numbers.
\end{CornellOverviewBox}\begin{CornellSummaryBox}{Summary}
{\sffamily\bfseries\color{Accent} Summary}\par\vspace{3pt}Section 5.9 focuses on Preprocessing numbers. Use the listed grammar productions together with the semantic constraints and disambiguation rules referenced by the section. When applying it, preserve the distinction between mandatory requirements, recommendations, permissions, and behavior deliberately left undefined, unspecified, or implementation-defined.
\end{CornellSummaryBox}

\end{document}
